\section{Diagnostic Utility}
\label{sec:diagnostic}
\begin{figure}[t]
\centering
\includegraphics[width=0.8\columnwidth]{Figures/raw_figures/figure3.pdf}
\vspace{-0.15in}
\caption{
\textbf{Failure composition on the A1--A2--A3 ladder.}
Incorrect closed-book responses by failure type across three Korean LLMs.
A1 corresponds to the seed. \textit{Other} includes invalid outputs and defective questions.
}
\label{fig:ladder}
\vspace{-0.15in}
\end{figure}
\begin{figure*}[t]
\centering
\includegraphics[width=\textwidth]{Figures/raw_figures/figure4.pdf}
\vspace{-0.25in}
\caption{
\textbf{Effects of knowledge and application interventions.}
Accuracy changes across closed-book, knowledge (+K), and knowledge and application (+K+A) conditions for A2 (78 questions) and A3 (209 questions).
Each bar shows closed-book accuracy, the additional gains from +K and +K+A, and the remaining error rate.
}
\label{fig:intervention}
\vspace{-0.1in}
\end{figure*}

Finally, we examine whether the contrasts built by \name{} can separate models' failures to recall knowledge from their failures to apply it.


\subsection{Experimental Setup}
\label{sec:diagnostic_setup}

We evaluate eight models\footnote{Reasoning is enabled for all models except Aya Expanse 32B.} on the 287 final questions from Sec.~\ref{sec:construction_quality}, covering four groups.
\begin{itemize}[leftmargin=3.5mm, topsep=2pt, itemsep=2pt, parsep=0pt]
    \item[$\circ$] \textit{Korean-specific}: K-EXAONE~\citep{choi2026k}, A.X K1~\citep{cheon2026ax}, and Solar Open~\citep{park2026solar}, as in Sec.~\ref{sec:failure}.
    \item[$\circ$] \textit{Multilingual-focused}: Aya Expanse 32B~\citep{dang2024aya}.
    \item[$\circ$] \textit{Open-weight frontier}: GLM-4.7~\citep{zai2025glm47} and Qwen3.5-397B-A17B~\citep{qwen2026qwen35}.
    \item[$\circ$] \textit{Proprietary frontier}: Claude Sonnet 5.5~\citep{anthropic2026sonnet55} and GPT-6.1 Sol~\citep{openai2026gpt61sol}.
\end{itemize}

Each question is answered under three conditions: \textit{closed-book}, with only the question and answer options; \textit{knowledge} (+K), with three facts drawn from the knowledge units of the anchor bundle; and \textit{knowledge and application} (+K+A), with further statements drawn from the relations of the bundle.
Both additions are written without access to the question and are shared by all questions built from the same seed.

For matched comparisons, we use the 30 seeds that have both an A2 and an A3 contrast.
All of the seeds are A1, so together with the seed question they form a complete A1--A2--A3 ladder.
%We label the incorrect closed-book responses of the three Korean LLMs with the procedure of Sec.~\ref{sec:failure}, using the solution written during construction in place of silver reasoning.

\subsection{Performance across Application Levels}
\label{sec:diagnostic_performance}

Table~\ref{tab:matched} reports closed-book accuracy across application levels.
Across all questions, seven of the eight models achieve lower accuracy on A3 than on A2.
Since most questions at the two levels come from different seeds, we also compare them on the 30 seeds that have both an A2 and an A3 contrast.
On the matched A1--A2--A3 ladders, where the underlying knowledge is held fixed, all eight models perform no better on A3 than on A2.
This consistent pattern suggests that increasing application demand makes questions more challenging even when the required knowledge is controlled.

\subsection{Failure Diagnosis and Intervention}
\label{sec:diagnostic_failure}

\paragraph{Failure profiles.}
On the matched A1--A2--A3 ladders, the share of application failures of the three Korean LLMs rises from 37.3\% at A1 to 60.0\% at A2 and 76.9\% at A3, while the share of knowledge failures falls from 41.8\% to 14.3\% (Figure~\ref{fig:ladder}).
The same shift appears across all 165 seeds, where application failures account for 32.5--46.6\% of incorrect responses on the seeds and 67.2--78.8\% on their contrasts.
This suggests that when the required knowledge is controlled, increasing application demand shifts failures from knowledge to application.

\paragraph{Knowledge and application intervention.}
Figure~\ref{fig:intervention} shows how accuracy changes as knowledge and application statements are provided.
Providing knowledge resolves most errors at A2, where accuracy reaches 83.3--100\% for all models except Aya Expanse 32B, but leaves more errors unresolved at A3, where accuracy remains at 62.7--94.3\%.
Adding the application statements then improves A3 accuracy for all eight models, by 2.9 to 15.3 points, with smaller gains at A2.
These results suggest that knowledge alone does not suffice when application demand is high, and that many of the remaining errors lie in applying the knowledge rather than in lacking it.